\documentclass[10pt,a4paper]{amsart}
\usepackage[margin=19mm]{geometry}
\usepackage{amsmath,amssymb,mathtools}
\usepackage[scr=rsfs,cal=euler]{mathalfa}
\usepackage{xfrac}
\usepackage[unicode,hidelinks,bookmarks=false]{hyperref}
\newcommand{\ie}{i.e.\ }
\newcommand{\cf}{cf.\ }
\newcommand{\resp}{respectively\ }
\newcommand{\Subsection}{\subsection}
\providecommand{\nobreakdash}{\nobreak\hbox{-}}
\setlength{\emergencystretch}{2em}
\multlinegap0pt
\usepackage{amssymb}
\newtheorem{theorem}{Theorem}
\newtheorem{lemma}{Lemma}
\DeclareMathOperator{\divOp}{div}
\DeclareMathOperator{\modOp}{mod}
\DeclareMathOperator{\ord}{ord}
\title{$p$-adic $\lambda$ functions for cyclic Mumford curves}
\author{Yaacov Kopeliovich}
\date{Source: arXiv:2512.23913v1 (CC BY 4.0)}
\begin{document}
\maketitle

\noindent\emph{Source and licence.} $p$-adic $\lambda$ functions for cyclic Mumford curves. Version arXiv:2512.23913v1 (CC BY 4.0), distributed by arXiv under the Creative Commons Attribution 4.0 International licence, http://creativecommons.org/licenses/by/4.0/, as recorded in the arXiv licence metadata for this exact version and shown on the abstract page https://arxiv.org/abs/2512.23913v1. Full licence notice: \texttt{licenses/arxiv-2512.23913-CC-BY-4.0.txt}.\par\smallskip

\noindent\textit{Editorial scope.} Counted unit: the article's theorem dimensionformula (source0.tex lines 140--152). The article's complete original proof of it is delivered with it (source0.tex lines 153--228, one \texttt{proof} environment of 76 lines). No mathematics is written, repaired or re-derived: the statement and the proof are the article's own lines, in the article's own notation, and no retained line is substituted or re-flowed.  No further source-local context is needed: every cross-reference of every retained line resolves inside the two delivered spans.  Nothing was omitted from any retained span, so the package carries no omission record.  No retained line contains a citation, so the wrapper carries no bibliography and no bibliography entry.  No retained line contains a figure environment, an image-inclusion command or drawing code, so no image is delivered and the package declares no asset dependency.  Editorial formatting only, in addition: an AMS \texttt{amsart} preamble that restates the article's own declarations and macros used by the retained lines, namely the environments \texttt{theorem}, \texttt{lemma} on separate counters, together with the standard packages those lines need.  The retained environments print in this document as Theorem 1, Lemma 1 in order of appearance, whereas the article prints the same items with the numbering of its own full document.  The complete native source of the article as distributed in the arXiv e-print for this exact version is bundled unchanged as \texttt{sources/191839/source0.tex}.
\begin{theorem}\label{dimensionformula}
Let $D\,{=}\,\sum_{i=1}^n d_i D_i$ be a divisor supported on finite branch points, 
and $d_j \,{\in}\, \{0,\dots,p-1\}$. 
Define the set  $\{g_k\}_{k=0}^{p-1}$ by the equations 
\begin{equation}
\frac{1}{p}\bigg(\deg D-\sum_{j=1}^n \overline{k r_j + d_j}
+\overline{\sum_{j=1}^n k r_j}\bigg) +1=g_k,\quad   0\leq k\leq p-1,
\end{equation}
where $\overline{r}$ denotes the remainder of the division of $r$ by $p$. Then 
\begin{equation}
\dim H^{0}(X,\mathcal{O}(-D))=\sum_{k=0}^{p-1} \max(0,g_k).
\end{equation}
\end{theorem}

\begin{proof}
The curve $X$ has a cyclic automorphism, given explicitly through 
$$T:(x,y)\mapsto (x,\zeta y),$$ 
where $\zeta$ is a $p$-th root of unity. 
$D$ is invariant under the action of $T$. Thus,  $T$ acts on $H^{0}(X,\mathcal{O}(-D))$.
We conclude that 
$$H^{0}(X,\mathcal{O}(-D)) \,{=}\,\bigoplus_{\chi\in \mathbb{Z}_p^{\ast}} V_{\chi},$$
where $V_{\chi}$ is the eigenspace of $T$ with a character 
$\chi\,{:}\,{\mathbb{Z}_p} \,{\to}\, \mathbb{K}$. 
Let $n_{\chi}\,{=}\,\dim V_{\chi}$. Clearly, 
$$\dim H^{0}(X,\mathcal{O}(-D)) = \sum_{\chi\in \mathbb{Z}_p^{\ast}} n_{\chi}.$$ 
We  find $n_{\chi}$ for all $\chi$.

 The characters $\chi$ of $T$ are of the form $\zeta^{\ell}$ for some $\ell$,
$0\,{\leq}\, \ell \,{\leq}\, p\,{-}\,1$. Hence, each $\chi$ is assigned with 
$\ell$ such that $\forall f \,{\in}\, V_{\chi}$ $f/y^{\ell}$ has a trivial action of $T$ on it. 
Thus, $f/y^{\ell}$ is a pullback of a function from $\mathbb{P}^{1}(\mathbb{K})$.
Let $H_{\ell}^{0}$ be the subspace of $H^0(X,\mathcal{O}(-D-\divOp(y^{\ell})))$ 
composed of all pullbacks from functions on $\mathbb{P}^{1}(\mathbb{K}).$
\begin{lemma}
$\exists D_{\ell}$, a divisor on $\mathbb{P}^{1}(\mathbb{K})$, 
such that $H^{0}(\mathbb{P}^{1}(\mathbb{K}),\mathcal{O}(-D_{\ell}))$ is isomorphic to $H^{0}_{\ell}$.
\end{lemma}
\begin{proof}
If $B_j$ is a ramification point appearing in the divisor $D$
with multiplicity $d_j$, 
let $m_j\,{=}\,[(\ell r_j\,{+}\,d_j)/p ]\,{\times}\, p$, that is,  $m_j$ is the maximal number such that 
$m_j \,{\leq}\, \ell r_j\,{+}\,d_j$ and $m_j \,{=}\, 0 \modOp p$. 
Let  $Q_{a_j}$ represent the point that corresponds 
to $a_j$ in $\mathbb{P}^{1}(\mathbb{K})$. Define
\begin{equation}
D_{\ell} = \sum_{j=1}^n \frac{m_j}{p} Q_{\lambda_j}-
\bigg[\ell \frac{\sum_{j=1}^n r_j}{p}\bigg] \infty.
\end{equation} 
We show that $D_{\ell}$ is the desired divisor. Let $h:\mathbb{P}^{1}(\mathbb{K})\to \mathbb{P}^{1}(\mathbb{K})$ 
such that $\divOp (h) \,{\geq}\, {-} D_{\ell}$.  By $\widehat{h}$ we denote a lift of $h$, 
and so $\widehat{h} \,{\in}\, H_{\ell}^{0}$.  Then
\begin{equation}
-\divOp (\widehat{h})\geq \sum_{j=1}^n -\tilde{d}_j B_j+ \ell \sum_{j=1}^n  r_i\infty.
\end{equation}
We require $\tilde{d}_j \,{\leq}\, \ell r_j \,{+}\, d_j$. 

Conversely, if $\widehat{h} \in H_{\ell}^{0}$ and $\widehat{h}$ is a lift of $h$, 
then by definition $$\ord_{P_j} \widehat{h} \geq - (\ell r_j + d_j).$$ 
We  have  $\ord_{P_j} \widehat{h}  \,{=}\,0\modOp p$,
since $\widehat{h}$ is a lift of $h$. 
We conclude, that $\ord_{Q_j} h \,{\geq}\, m_j/p$, 
and hence, $h\in H_{\ell}^{0}$, as required. 
\end{proof}
Now we compute $\deg D_{\ell}$: 
\begin{equation}
\deg D_{\ell} =\frac{1}{p}\bigg(\sum_{j=1}^{n} m_j \bigg)
-\bigg[\frac{\sum_{j=1}^n \ell r_j}{p}\bigg]. 
\end{equation}
By definition, $m_j\,{=}\, \ell r_j \,{+}\, d_j \,{-}\,\overline{\ell r_j \,{+}\, d_j}$. Hence, we can write
\begin{equation}
\begin{split}
\deg D_{\ell} &= \frac{1}{p}\bigg(\sum_{j=1}^n (\ell r_j + d_j)
-\sum_{j=1}^n \overline{\ell r_j + d_j}-\sum_{j=1}^n \ell r_j+\overline{\sum_{j=1}^n \ell r_j}\bigg) \\
&=\frac{1}{p}\bigg(\deg D -\sum_{j=1}^n\overline{\ell r_j + d_j}+\overline{\sum_{j=1}^n \ell r_j}\bigg).
\end{split}
\end{equation}
Applying the Riemann-Roch theorem to $D_{\ell}$, we obtain
$$\dim H^0 \big(\mathbb{P}^{1}(\mathbb{K}),\mathcal{O}(-D_{\ell}) \big)
= \max (\deg D_{\ell} +1,0). $$ 
Therefore,
\begin{equation}
\begin{split}
\dim H^0 \big(X,\mathcal O(-D)\big) &= \sum_{\chi_{\ell} \in Z_p^{*}} \max(\deg D_{\ell} +1,0)\\
&=\sum_{\ell=0}^{p-1} \max\bigg(\frac{1}{p}\bigg(\deg D - \sum_{j=1}^n \overline{\ell r_j + d_j}
+\overline{\sum_{j=1}^n \ell r_j}\bigg)+1,0\bigg).
\end{split}
\end{equation}
The last sum equals 
$\sum_{\ell=0}^{p-1} \max(0,g_{\ell})$, which completes the proof of the theorem. 
\end{proof}

\end{document}
